\section{Einleitung}
\subsection{Motivation und Kontext}
\subsection{Problemstellung}



\subsubsection{Derzeitige Forschungslücken}

Aktuelle Forschungen zum Einsatz von \gls{ids}, wie die von \cite{kladz2025ids}, 
fokussieren sich primär auf die reine Validierung von Modellen. 
Eine Lücke besteht in der aktiven Ergänzung fehlender oder 
unvollständiger Informationen in das Modell, insbesondere in der Abgebotsphase von Bauprojekten.
 
Wie \cite{de2024enriching} betonen, besitzt \gls{ids} das Potenzial, über die reine Validierung hinaus 
zur Modellanreicherung beizutragen. Konkrete Implementierungen, die speziell auf die Anforderungen von 
Bauunternehmen in der Angebotsphase zugeschnitten sind, fehlen.

Diese Arbeit setzt hier an, indem sie aufzeigt, wie offene Standards (\gls{ids}, \gls{ifc}) 
und Open-Source-Technologien (IfcOpenShell, Blender) kombiniert werden können, um hier einen 
durchgängigen \gls{bim}-Prozess zu realisieren.


\subsection{Zielsetzung und Forschungsfrage}
\subsection{Aufbau der Arbeit}